\documentclass[10.5pt]{article}
\usepackage[margin=0.85in, top=0.75in, bottom=0.75in]{geometry}
\usepackage{microtype}
\usepackage{parskip}
\usepackage[colorlinks=true,linkcolor=black,citecolor=black,urlcolor=blue]{hyperref}

\title{\textbf{Grounded Reasoning: Why Embedding an LLM\\[4pt]
in an Executable Environment Improves Performance}}
\author{}
\date{May 2026}

\begin{document}
\maketitle
\vspace{-1em}

\subsection*{The Problem with Pure Reasoning}

Large language models (LLMs) are trained to predict plausible continuations of
text. This produces models that are remarkably fluent and knowledgeable, but it
carries a fundamental limitation: the model's ``beliefs'' about the world are
encoded as distributional biases in its weights, and there is no mechanism by
which incorrect reasoning is corrected by contact with reality. When an LLM
reasons about code purely in natural language---thinking about what a program
\emph{will} do rather than running it---it generates whatever continuation is
most likely under its training distribution. That distribution reflects the
frequency with which descriptions appear in human text, not actual execution
outcomes. The result is a systematic tendency toward plausible-sounding but
incorrect reasoning: the model cannot reliably distinguish correct programs from
subtly broken ones, because both appear in human-written discussion with similar
surface characteristics.

\subsection*{Empirical Evidence for Execution Loops}

The practical superiority of embedding LLMs in execution loops is now
well-documented. Program-Aided Language Models (PAL) \cite{gao2023pal}
demonstrated that directing an LLM to express reasoning steps as executable
Python---and then running that code---substantially outperforms chain-of-thought
prompting on mathematical and symbolic tasks. The mechanism is transparent:
arithmetic and logical operations are delegated to the interpreter, which
executes them exactly rather than approximately. ReAct \cite{yao2023react}
formalised this as a general framework, interleaving language reasoning traces
with actions in an environment (including code execution). Crucially, ReAct
outperformed both pure reasoning (no actions) and pure acting (no reasoning),
suggesting the value lies in the combination of hypothesis formation and
empirical testing. Reflexion \cite{shinn2023reflexion} extended this further:
agents that receive execution feedback and are prompted to reflect verbally on
failures before retrying substantially outperform single-shot generation on
programming benchmarks, constituting a lightweight form of iterative hypothesis
refinement.

SWE-bench \cite{jimenez2024swebench}, a benchmark of real GitHub issues
requiring code modification and test-suite verification, demonstrated that the
gap between single-pass generation and iterative agents with execution feedback
is large and practically significant. InterCode \cite{yang2023intercode}
treated code generation as a partially observable Markov decision process, with
execution outcomes as observations, and found consistent gains across multiple
programming environments.

\subsection*{Why Does It Work? Partial Theoretical Accounts}

The empirical picture is clear; the theoretical understanding is less settled.
Several complementary accounts have been offered, none fully satisfactory.

\textbf{Grounding via unambiguous feedback.} The most direct account is that
execution provides feedback the model cannot hallucinate. A Python traceback is
a deterministic output of a formal system---categorically different from the
model's self-evaluation, which is itself a text prediction subject to the same
distributional biases that produced the error. The environment acts as an oracle
external to the model's uncertainty.

\textbf{Offloading formal computation.} PAL and related work suggest that LLMs
are unreliable at tasks requiring exact computation---not from lack of relevant
knowledge, but because approximating exact arithmetic through token prediction is
inherently lossy. Generating code that delegates computation to the interpreter
avoids attempting what the model does poorly. This mirrors what humans do when
using a calculator rather than performing mental arithmetic for critical
computations.

\textbf{Error decomposition and localisation.} In multi-step reasoning, errors
compound: a single incorrect intermediate result propagates through subsequent
steps. Execution grounds each step. If a program produces a wrong intermediate
value, the test suite or downstream execution reveals this before compounding
occurs. The loop converts a problem of unbounded error propagation into one of
bounded, localised error correction.

\subsection*{On the Reservoir Computing Analogy}

One might ask whether the programming runtime plays a role analogous to a
reservoir in reservoir computing \cite{jaeger2001,maass2002}: a fixed, complex
dynamical system that transforms inputs into a high-dimensional state space,
with the LLM acting as a learned readout. The analogy is suggestive but should
be treated with caution. In classical reservoir computing, the reservoir is a
random, typically chaotic network whose high-dimensional nonlinear dynamics are
exploited \emph{precisely because} they are unlearned and fixed. The programming
runtime, by contrast, is structured, deterministic, and semantically
interpretable. More importantly, the LLM does not merely read from the runtime's
state; it generates instructions that determine what the runtime computes. The
control flow runs in the opposite direction from standard reservoir computing.

A more defensible analogy is the extended mind hypothesis \cite{clark1998}: that
cognitive processes can be distributed between an agent and its environment when
the environment reliably stores and transforms information. On this view, the
executing program is part of the reasoning process, not merely a tool external
to it. Whether this framing is scientifically productive or merely a useful
metaphor is itself an open question.

\subsection*{Honest Assessment and Open Questions}

The core empirical finding---that LLMs in execution loops substantially
outperform single-pass generation on coding and formal reasoning tasks---is
robust and reproducible. The theoretical explanation is less settled. The
grounding and offloading accounts are plausible and consistent with the data,
but they do not constitute a mechanistic theory of why language model weights
fail to reliably simulate formal computation internally. It remains unclear
whether this is a fundamental limitation of autoregressive prediction as a
training objective, a consequence of insufficient scale, or an architectural
constraint that future models might overcome.

It is also worth being honest about what execution loops \emph{cannot} do: they
do not improve reasoning about problems that cannot be formalised or tested.
The benefit is specific to domains with deterministic, observable outcomes.
Finally, the question of whether empirical exploration in a formal environment
is qualitatively different from internal reasoning---or merely more
reliable---remains philosophically and empirically open. The honest summary is:
we have strong evidence that it works, a plausible but incomplete story about
why, and a significant theoretical gap remaining.

\begin{thebibliography}{9}
\bibitem{gao2023pal} Gao, L.\ et al.\ PAL: Program-aided language models.
  \textit{Proc.\ ICML}, PMLR 202:10764--10799, 2023.

\bibitem{yao2023react} Yao, S.\ et al.\ ReAct: Synergizing reasoning and acting
  in language models. \textit{ICLR}, 2023.

\bibitem{shinn2023reflexion} Shinn, N.\ et al.\ Reflexion: Language agents with
  verbal reinforcement learning. \textit{NeurIPS}, 2023.

\bibitem{jimenez2024swebench} Jimenez, C.\,E.\ et al.\ SWE-bench: Can language
  models resolve real-world GitHub issues? \textit{ICLR}, 2024.

\bibitem{yang2023intercode} Yang, J.\ et al.\ InterCode: Standardizing and
  benchmarking interactive coding with execution feedback. \textit{NeurIPS}, 2023.

\bibitem{jaeger2001} Jaeger, H.\ The ``echo state'' approach to analysing and
  training recurrent neural networks. \textit{GMD Report} 148, 2001.

\bibitem{maass2002} Maass, W., Natschl\"{a}ger, T., Markram, H.\ Real-time
  computing without stable states: A new framework for neural computation based
  on perturbations. \textit{Neural Computation} 14(11):2531--2560, 2002.

\bibitem{clark1998} Clark, A., Chalmers, D.\ The extended mind.
  \textit{Analysis} 58(1):7--19, 1998.
\end{thebibliography}

\end{document}
